\documentclass[11pt]{article}
\usepackage[utf8]{inputenc}
\usepackage{amsmath,amssymb}
\usepackage{hyperref}
\title{Efficient Protein Structure Prediction Deep Learning: A Resource-Aware Approach}
\author{Anonymous}
\date{}
\begin{document}
\maketitle

\begin{abstract}
This paper studies Protein Structure Prediction Deep Learning. Grounded in a real, citable literature \cite{ref1,ref2,ref3,ref4,ref5,ref6,ref7,ref8},
we propose a focused method and report \textbf{illustrative} experiments.
All references below are real and verifiable (OpenAlex/arXiv); the reported
numbers are simulated to illustrate intended behaviour.
\end{abstract}

\section{Introduction}
Protein Structure Prediction Deep Learning is an active research area. This work targets a concrete gap identified from
the recent literature and contributes a method together with an evaluation protocol.

\section{Related Work (real literature)}
The following works, retrieved from public bibliographic APIs, frame our study:
\begin{itemize}
\item Jumper et al. (2021) — "Highly accurate protein structure prediction with AlphaFold", Nature (cited 46082).
\item Abramson et al. (2024) — "Accurate structure prediction of biomolecular interactions with AlphaFold 3", Nature (cited 14666).
\item Senior et al. (2020) — "Improved protein structure prediction using potentials from deep learning", Nature (cited 3565).
\item Rives et al. (2021) — "Biological structure and function emerge from scaling unsupervised learning to 250 million protein sequences", Proceedings of the National Academy of Sciences (cited 3246).
\item Watson et al. (2023) — "De novo design of protein structure and function with RFdiffusion", Nature (cited 2075).
\item Dauparas et al. (2022) — "Robust deep learning–based protein sequence design using ProteinMPNN", Science (cited 1946).
\end{itemize}

\section{Method}
We reduce the dominant computational cost via adaptive resource allocation, activating only the components needed per input. We instantiate this idea for Protein Structure Prediction Deep Learning and analyse its cost/quality trade-off.

\section{Experiments (Illustrative)}
\begin{center}
\begin{tabular}{lcc}
System & Primary Metric & Efficiency \\ \hline
Strong baseline & 0.812 & 1.00$\times$ \\
Ablation & 0.828 & 0.92$\times$ \\
\textbf{Ours} & \textbf{0.851} & \textbf{0.71$\times$}
\end{tabular}
\end{center}
\emph{Metrics simulated to illustrate the intended effect; not measured.}

\section{Conclusion}
We connected a real literature to a concrete method for Protein Structure Prediction Deep Learning. Future work will
validate the illustrative results with real experiments.

\begin{thebibliography}{9}
\bibitem{ref1} John Jumper, Richard Evans, Alexander Pritzel et al.. Highly accurate protein structure prediction with AlphaFold. \emph{Nature}, 2021. DOI:10.1038/s41586-021-03819-2
\bibitem{ref2} Josh Abramson, Jonas Adler, Jack Dunger et al.. Accurate structure prediction of biomolecular interactions with AlphaFold 3. \emph{Nature}, 2024. DOI:10.1038/s41586-024-07487-w
\bibitem{ref3} Andrew Senior, Richard Evans, John Jumper et al.. Improved protein structure prediction using potentials from deep learning. \emph{Nature}, 2020. DOI:10.1038/s41586-019-1923-7
\bibitem{ref4} Alexander Rives, Joshua Meier, Tom Sercu et al.. Biological structure and function emerge from scaling unsupervised learning to 250 million protein sequences. \emph{Proceedings of the National Academy of Sciences}, 2021. DOI:10.1073/pnas.2016239118
\bibitem{ref5} Joseph L. Watson, David Juergens, Nathaniel R. Bennett et al.. De novo design of protein structure and function with RFdiffusion. \emph{Nature}, 2023. DOI:10.1038/s41586-023-06415-8
\bibitem{ref6} Justas Dauparas, Ivan Anishchenko, Nathaniel R. Bennett et al.. Robust deep learning–based protein sequence design using ProteinMPNN. \emph{Science}, 2022. DOI:10.1126/science.add2187
\bibitem{ref7} Hakime Öztürk, Arzucan Özgür, Elif Özkırımlı. DeepDTA: deep drug–target binding affinity prediction. \emph{Bioinformatics}, 2018. DOI:10.1093/bioinformatics/bty593
\bibitem{ref8} Jianyi Yang, Ivan Anishchenko, Hahnbeom Park et al.. Improved protein structure prediction using predicted interresidue orientations. \emph{Proceedings of the National Academy of Sciences}, 2020. DOI:10.1073/pnas.1914677117
\end{thebibliography}
\end{document}
